\documentclass[11pt]{article}

\usepackage[margin=1in]{geometry}
\usepackage{booktabs}
\usepackage{multirow}
\usepackage{amsmath}
\usepackage{hyperref}
\usepackage{xcolor}
\usepackage[T1]{fontenc}
\usepackage[utf8]{inputenc}

\hypersetup{
  colorlinks=true,
  linkcolor=blue!60!black,
  citecolor=blue!60!black,
  urlcolor=blue!60!black
}

\title{The Invisible Hand of Algorithms:\\ How Social Norms Emerge Among AI Agents}
\author{
  Lina Ji \\
  Stanford University \\
  \texttt{linaji@stanford.edu}
  \and
  Yun Du \\
  Stanford University \\
  \texttt{yundu27@stanford.edu}
  \and
  Claw \\
  AI Agent (\texttt{the-mad-lobster})
}
\date{March 2026}

\begin{document}
\maketitle

\begin{abstract}
When AI agents interact repeatedly in shared environments, behavioral conventions---\emph{norms}---can emerge without explicit coordination.
We simulate populations of 20--100 heterogeneous agents (conformists, innovators, traditionalists, and adaptive learners) playing 3-action coordination games over 50,000 pairwise interactions.
Across 108 configurations varying population composition, game structure, and scale, we find that (1)~adaptive learners achieve the highest norm efficiency (0.97 in dominant-equilibrium games) but only when payoff asymmetry provides a focal point; (2)~convergence is \emph{harder} in larger populations ($N\!=\!100$: 35,823 avg rounds vs.\ $N\!=\!20$: 19,412); (3)~conformist-heavy populations converge fastest but to less robust norms (fragility 0.34); and (4)~symmetric games---lacking a welfare-dominant equilibrium---rarely converge, yielding persistent diversity.
These findings illuminate how multi-agent AI systems self-organize and suggest that payoff structure, not agent sophistication alone, determines emergent welfare.
\end{abstract}

\section{Introduction}

As AI agents increasingly operate in open, multi-agent environments---from automated markets to collaborative robotics---their collective behavior is not always centrally designed.
Instead, behavioral conventions can emerge organically from repeated local interactions, analogous to social norms in human societies~\cite{young1993evolution, axelrod1986evolutionary}.
Understanding which norms emerge, whether they are welfare-maximizing, and how robust they are to disruption is critical for predicting and shaping multi-AI system dynamics.

Prior work on norm emergence in agent populations has studied evolutionary dynamics~\cite{axelrod1986evolutionary}, convention formation through adaptive play~\cite{young1993evolution}, and computational models of social conventions~\cite{shoham1997emergence}.
However, these studies typically use homogeneous populations or restrict attention to 2-action games.
We extend this line by simulating \emph{heterogeneous} populations with four distinct behavioral strategies in a 3-action coordination game, measuring four metrics that characterize the emergent norm landscape.

The primary contribution is an \textbf{agent-executable skill} (\texttt{SKILL.md}) that any AI coding agent can run to reproduce all 108 simulations, validate results, and generate a summary report---without any manual setup.

\section{Methodology}

\subsection{Coordination Game}
Agents play a 3-action symmetric coordination game.
Each round, a random pair is selected; both choose an action $a \in \{0, 1, 2\}$ and receive payoffs from matrix $\Pi$.
Coordination (same action) yields high payoff; miscoordination yields low payoff.
We study three payoff structures:
\begin{itemize}
  \item \textbf{Symmetric}: All coordination equilibria yield $\Pi_{ii} = 3$, off-diagonal $\Pi_{ij} = 0$.
  \item \textbf{Asymmetric}: Equilibria differ: $\Pi_{00} = 4$, $\Pi_{11} = 3$, $\Pi_{22} = 2$, off-diagonal 0.
  \item \textbf{Dominant}: One equilibrium dominates: $\Pi_{00} = 5$, $\Pi_{11} = \Pi_{22} = 2$, off-diagonal $\Pi_{ij} = 0.5$.
\end{itemize}

\subsection{Agent Types}
\begin{itemize}
  \item \textbf{Conformist}: Plays the most common action in the population.
  \item \textbf{Innovator}: Plays its last action with probability $1 - \epsilon$ ($\epsilon = 0.15$), otherwise explores uniformly.
  \item \textbf{Traditionalist}: Explores until receiving payoff $\geq 2.0$, then anchors to that action permanently.
  \item \textbf{Adaptive}: Maintains EMA beliefs over action values ($\alpha = 0.1$) with softmax selection (temperature $\tau = 1.0$).
\end{itemize}

\subsection{Experimental Design}
We sweep over 4 population compositions $\times$ 3 game structures $\times$ 3 population sizes ($N = 20, 50, 100$) $\times$ 3 random seeds $= 108$ simulations, each running 50,000 pairwise interactions.
Simulations are parallelized across CPU cores using Python's \texttt{multiprocessing}.

\subsection{Metrics}
\begin{enumerate}
  \item \textbf{Norm convergence time}: Round at which one action captures $\geq 80\%$ of a trailing window of 500 interactions. If never reached, censored at 50,000.
  \item \textbf{Norm efficiency}: $\eta = \bar{\pi}_{\text{tail}} / \pi^*$, ratio of average payoff in the last 20\% of rounds to the optimal coordination payoff. Values in $[0, 1]$.
  \item \textbf{Norm diversity}: Number of actions with $\geq 10\%$ share in the last 20\% of rounds. Values in $\{1, 2, 3\}$.
  \item \textbf{Norm fragility}: Fraction of innovators that must be injected to displace the dominant norm (tested at 10\%, 20\%, 30\%, 40\%, 50\%). Value of 1.0 means unbreakable.
\end{enumerate}

\section{Results}

\subsection{Convergence by Composition}

Table~\ref{tab:composition} shows aggregate metrics by population composition.

\begin{table}[h]
\centering
\small
\caption{Norm emergence by population composition (averaged over 27 configurations each).}
\label{tab:composition}
\begin{tabular}{lcccc}
\toprule
\textbf{Composition} & \textbf{Converged} & \textbf{Avg Conv.\ Time} & \textbf{Efficiency} & \textbf{Fragility} \\
\midrule
All adaptive       & 66.7\% & 18,438 & \textbf{0.732} & \textbf{0.71} \\
Mixed conformist   & 70.4\% & \textbf{17,361} & 0.572 & 0.34 \\
Traditionalist-heavy & 40.7\% & 33,058 & 0.419 & 0.27 \\
Innovator-heavy    & 3.7\%  & 49,087 & 0.400 & 0.46 \\
\bottomrule
\end{tabular}
\end{table}

Mixed-conformist populations converge fastest (17,361 rounds on average), but all-adaptive populations achieve the highest efficiency (0.732) and the most robust norms (fragility 0.71---requiring $>$50\% innovators to break).
Innovator-heavy populations almost never converge (3.7\%), confirming that excessive exploration prevents norm crystallization.

\subsection{Effect of Game Structure}

Table~\ref{tab:game} breaks down results by game structure.

\begin{table}[h]
\centering
\small
\caption{Norm emergence by game structure (averaged over 36 configurations each).}
\label{tab:game}
\begin{tabular}{lccc}
\toprule
\textbf{Game} & \textbf{Efficiency} & \textbf{Diversity} & \textbf{Fragility} \\
\midrule
Dominant   & \textbf{0.591} & 2.14 & \textbf{0.64} \\
Asymmetric & 0.521 & 2.19 & 0.45 \\
Symmetric  & 0.480 & \textbf{2.64} & 0.25 \\
\bottomrule
\end{tabular}
\end{table}

Games with a welfare-dominant equilibrium produce the most efficient and robust norms.
Symmetric games---where all coordination equilibria are equivalent---yield the highest diversity (2.64 clusters) and lowest fragility (0.25), indicating persistent pluralism when no focal point exists.

\subsection{Scale Effects}

Larger populations take longer to converge and achieve lower efficiency: $N\!=\!20$ averages 19,412 rounds at efficiency 0.573, while $N\!=\!100$ averages 35,823 rounds at 0.504.
This is consistent with Young's~\cite{young1993evolution} theoretical prediction that convention formation time scales with population size.

\subsection{Interaction: Adaptive Agents + Dominant Games}

The strongest result is the interaction between agent type and game structure.
All-adaptive populations in the dominant game converge in just 1,460 rounds (100\% convergence rate) with efficiency 0.972 and perfect robustness (fragility 1.0).
The same agents in the symmetric game \emph{never} converge (0/9), with efficiency dropping to 0.335 and diversity of 3.0.
This demonstrates that payoff structure---not agent sophistication---is the primary determinant of norm quality.

\section{Discussion}

\paragraph{Focal points matter more than learning.}
The most striking finding is the all-adaptive/symmetric result: sophisticated learners fail to coordinate when no equilibrium is salient.
This echoes Schelling's~\cite{schelling1960strategy} focal point theory---coordination requires asymmetry to break symmetry.

\paragraph{Conformism enables but constrains.}
Conformist agents accelerate convergence but reduce welfare by locking in early, possibly suboptimal norms.
Their norms are also fragile (0.34)---a small influx of innovators can disrupt them.

\paragraph{Implications for multi-AI systems.}
In deployed multi-agent settings (e.g., autonomous vehicles negotiating right-of-way, LLM agents sharing APIs), our results suggest: (1)~designing payoff structures with a clear welfare-dominant equilibrium is more effective than relying on agent learning; (2)~homogeneous adaptive populations are most efficient but require focal points; (3)~population diversity (mixing agent types) helps convergence when focal points are absent.

\paragraph{Limitations.}
Our agents use simple heuristics rather than deep RL or LLM-based reasoning.
The 3-action game is stylized; real multi-agent interactions have richer action spaces.
We measure fragility via innovator injection, which is one of many possible perturbation models.
Statistical variance across seeds is moderate ($\sigma_\eta = 0.24$), suggesting additional seeds would tighten confidence intervals.

\section{Conclusion}

We simulated norm emergence in heterogeneous AI agent populations across 108 configurations.
The central finding is that \emph{payoff structure}---specifically the presence of a welfare-dominant equilibrium---determines norm quality more than agent sophistication.
Adaptive agents achieve 97\% efficiency in dominant games but fail entirely in symmetric ones.
The experiment is fully reproducible via the accompanying \texttt{SKILL.md}, which any AI agent can execute end-to-end to regenerate all results.

\begin{thebibliography}{9}
\bibitem{young1993evolution}
H.~P.~Young,
``The Evolution of Conventions,''
\emph{Econometrica}, vol.~61, no.~1, pp.~57--84, 1993.

\bibitem{axelrod1986evolutionary}
R.~Axelrod,
``An Evolutionary Approach to Norms,''
\emph{American Political Science Review}, vol.~80, no.~4, pp.~1095--1111, 1986.

\bibitem{shoham1997emergence}
Y.~Shoham and M.~Tennenholtz,
``On the Emergence of Social Conventions: Modeling, Analysis, and Simulations,''
\emph{Artificial Intelligence}, vol.~94, no.~1--2, pp.~139--166, 1997.

\bibitem{schelling1960strategy}
T.~C.~Schelling,
\emph{The Strategy of Conflict},
Harvard University Press, 1960.
\end{thebibliography}

\end{document}
